\documentclass[11pt,a4paper]{article}
\usepackage[utf8]{inputenc}
\usepackage{geometry}
\geometry{margin=1in}
\usepackage{fancyhdr}
\usepackage{hyperref}
\usepackage{enumitem}
\usepackage{xcolor}
\usepackage{listings}

% Setup custom header/footer styling
\pagestyle{fancy}
\fancyhf{}
\lhead{\textbf{Agentic Wokflows: N8N}}
\rhead{ sujeet.banerjee@gmail.com \\ 
sujeet.banerjee.professional@gmail.com }
\lfoot{Simple Regn. workflow}
\rfoot{\thepage}

\title{\vspace{-2em}
\Huge \textbf{Agentic Wokflows: N8N} \\
\vspace{1em}
\large \textit{Practice Set: Building a simple Registration Workflow} 
}
\author{}
\date{}

\begin{document}

\maketitle
\vspace{-4em}
\hrule
\vspace{1.5em}

\section*{Objective}
Design and implement an automated workflow in n8n to process student registrations. This exercise will test your ability to capture input data, apply conditional logic, and structure data output dynamically.
\section*{}
Also outline test-scenarios for the workflow.

\section*{Core Requirements}
Create a workflow that takes the following inputs for a registration record:
\begin{itemize}
    \item \textbf{Name:} (String)
    \item \textbf{Experience (in years):} (Number)
\end{itemize}

The workflow must process this input and append the following new fields to the output record:
\begin{enumerate}
    \item \textbf{\texttt{assigned\_track}}:\\ If Experience is $< 5$ years, set to \texttt{'A'}. Otherwise, set to \texttt{'B'}.
    \item \textbf{\texttt{category}}:\\ If Name is exactly \texttt{'Sujeet'}, set to \texttt{'Trainer'}. Otherwise, set to \texttt{'Student'}.
    \item \textbf{\texttt{time}}:\\ Record the exact time of the registration processing (current timestamp).
    \item \textbf{Trainer Override Logic}:\\ If the \texttt{category} is determined to be \texttt{'Trainer'}, the \texttt{assigned\_track} must be cleared (set to blank or \texttt{null}).
\end{enumerate}

\section*{Implementation Guidelines}
\begin{itemize}
    \item \textbf{Trigger:} You may use any trigger to start the workflow (e.g., a \texttt{Manual Trigger} or \texttt{Chat Trigger}).
    \item \textbf{Mock Data Setup:} You may use an \texttt{Edit Fields (Set)} node immediately after the trigger to supply hardcoded input records for testing purposes.
    \item \textbf{Nodes to Explore:} Consider using the \texttt{If}, \texttt{Switch}, and \texttt{Edit Fields (Set)} nodes to achieve the required conditional routing and data mutation visually.
\end{itemize}

\section*{Optional Bonus Challenges}
For students looking to push their skills further, attempt the following advanced tasks:
\begin{enumerate}[label=\textbf{Bonus \arabic*:}]
    \item \textbf{Database Persistence (Postgres):} \\
    Sign up for a free online tier of CockroachDB (which is PostgreSQL-compatible). Add a \texttt{Postgres} node at the end of your workflow to insert the fully processed registration records into an online database table.
    
    \item \textbf{Code Node Optimization (JavaScript):} \\
    Replace the visual conditional flows (\texttt{If}/\texttt{Switch} nodes) entirely by using a single \texttt{Code} node. Write custom JavaScript within the node to calculate and set the \texttt{assigned\_track}, \texttt{category}, and \texttt{time}, including the override logic for the Trainer.
\end{enumerate}

\vspace{2em}
\hrule
\vspace{1em}
% \noindent \textit{Instructor Note: Ensure you test your workflow with multiple mock records (e.g., a student with 3 years experience, a student with 6 years experience, and 'Sujeet' with any amount of experience) to verify all logical branches function correctly.}

\end{document}